\documentclass[11pt]{article}

\usepackage[T1]{fontenc}
\usepackage{lmodern}
\usepackage[a4paper,margin=1in]{geometry}
\usepackage{amsmath}
\usepackage{amssymb}
\usepackage{amsthm}
\usepackage{mathtools}
\usepackage{xcolor}
\usepackage{listings}
\usepackage{booktabs}
\usepackage{hyperref}
\hypersetup{hidelinks}
\usepackage{microtype}
\setlength{\emergencystretch}{3em}

% ---- shared macros (mirroring the main paper) --------------------------
\newcommand{\code}[1]{\textsf{\small #1}}
\newcommand{\tool}{\textsc{StepCheck}}
\newcommand{\asl}{\textsc{Asl}}
\newcommand{\dom}{\operatorname{dom}}
\newcommand{\Obj}{\mathsf{Obj}}
\newcommand{\narrow}{\mathit{narrow}}
\newcommand{\place}{\mathit{place}}
\newcommand{\eff}{\mathit{eff}}
\newcommand{\Env}{\mathcal{E}}
\newcommand{\jval}{\mathit{jval}}
\newcommand{\jout}{\mathit{jout}}
\newcommand{\bind}{\mathit{bind}}
\newcommand{\asgn}{\mathit{asgn}}
% double square brackets without depending on stmaryrd
\providecommand{\llbracket}{\mathopen{[\![}}
\providecommand{\rrbracket}{\mathclose{]\!]}}

% ---- colours / listings -----------------------------------------------
\definecolor{kw}{RGB}{0,86,179}
\definecolor{cmt}{RGB}{0,128,0}
\definecolor{bg}{RGB}{248,248,248}
\definecolor{rule}{RGB}{120,120,120}
\lstdefinestyle{pseudo}{
  basicstyle=\ttfamily\small, backgroundcolor=\color{bg},
  frame=single, framesep=4pt, rulecolor=\color{rule!40},
  breaklines=true, columns=fullflexible, keepspaces=true,
  showstringspaces=false, xleftmargin=4pt, xrightmargin=4pt,
  commentstyle=\color{cmt}\itshape, keywordstyle=\color{kw}\bfseries,
  morekeywords={for,each,if,then,else,return,while,in,and,or,not,fire,let,do,foreach},
}

% ---- theorem environments ---------------------------------------------
\theoremstyle{plain}
\newtheorem{theorem}{Theorem}
\newtheorem{lemma}{Lemma}
\newtheorem{corollary}{Corollary}
\theoremstyle{definition}
\newtheorem{definition}{Definition}
\newtheorem{assumption}{Assumption}
\theoremstyle{remark}
\newtheorem{remark}{Remark}

\title{\textbf{Static Verification of AWS Step Functions with Sound
Data-Flow Analysis and Workflow Semantics}\\[4pt]\large Technical Report:
Formal Development and Proofs}
\author{Anonymous Author(s)}
\date{}

\begin{document}
\maketitle

\begin{abstract}
This technical report is the formal companion to the paper \emph{Static Verification of
AWS Step Functions with Sound Data-Flow Analysis and Workflow Semantics}~\cite{stepcheck},
which presents \tool, a static verifier for the Amazon States Language (\asl). The paper
states three guarantees and gives their intuition; here we give the complete development
behind them: the abstract interpretation underlying the sound data-flow provenance analysis
(domain, concretization, transfer functions, fixpoint, and the termination/convergence
guard), the workflow type discipline for contracts and typestate, and the downstream-aware
Saga-compensation check, with full proofs of all three theorems. Theorems~1--3 below are
numbered to match Theorems~1--3 of the main paper. We also give the \code{SC4011} algorithm
in full and note the conservative basis of the concurrency and temporal analyses. Section
and code references (e.g.\ \code{SC1101}) follow the main paper; this report is released with
the artifact and is self-contained with respect to the formal development, deferring only
the empirical material (the execution oracle and the evaluation) to the main paper.
\end{abstract}

\tableofcontents

\section{Scope and relation to the main paper}
\label{sec:scope}

The main paper keeps theorem statements and one-paragraph intuitions inline and defers their
proofs and the supporting formalism here, so the body stays within venue limits. This report
adds nothing to the claims; it discharges them.

We assume the \asl{} data plane as fixed in the paper's background~\cite{awsasl}. A workflow
is a state machine over eight state types (\texttt{Task}, \texttt{Choice}, \texttt{Parallel},
\texttt{Map}, \texttt{Wait}, \texttt{Pass}, \texttt{Succeed}, \texttt{Fail}). A single JSON
document is threaded between states and reshaped by the five-stage I/O pipeline
\texttt{InputPath}\,$\to$\,\texttt{Parameters}/\texttt{ItemSelector}\,$\to$\,%
\texttt{ResultSelector}\,$\to$\,\texttt{ResultPath}\,$\to$\,\texttt{OutputPath}. A
\texttt{.\$} field reference reads a JSONPath~\cite{jsonpath} into the current document, and
the runtime \emph{fails} the execution if such a reference resolves to an absent field.

The development below is faithful to the implementation (module \texttt{passes/dataflow.rs}
for Section~\ref{sec:tr:df}, \texttt{passes/contract.rs} and \texttt{passes/typestate.rs} for
Section~\ref{sec:tr:types}, and \texttt{passes/compensation.rs} for
Section~\ref{sec:tr:saga}); the execution oracle (\texttt{concrete.rs}) that empirically
witnesses Theorem~\ref{thm:soundness} is described in the main paper. Nested machines
(\texttt{Map} iterators and \texttt{Parallel} branches) are treated in
Remark~\ref{rmk:nested}.

Table~\ref{tab:tr:facts} records, for each fact a check needs, whether \asl{} exposes it
natively and where \tool{} obtains it otherwise (summarized in prose in the main paper's
background section).

\begin{table}[h]
\centering
\begin{tabular}{@{}p{4.6cm}p{2.4cm}p{5.6cm}@{}}
\toprule
\textbf{Fact a check needs} & \textbf{In \asl{}?} & \textbf{Source in \tool{}} \\
\midrule
Control flow and terminals & yes & native structural pass \\
JSONPath data movement & yes & native data-flow pass \\
Contracts and typestate & no & sidecar / DSL / schema metadata \\
Idempotency, persistence, compensator & no & declared premises; optional warning inference \\
Shared resources and time bounds & partial & conservative native passes \\
\bottomrule
\end{tabular}
\caption{Facts needed by \tool{} and where they come from.}
\label{tab:tr:facts}
\end{table}

\section{The data-flow abstract interpretation}
\label{sec:tr:df}

The provenance analysis is a forward abstract interpretation~\cite{cousot1977} over a lattice of
JSON \emph{document and variable shapes}: it over-approximates, at each state, the document and
workflow-variable store with which the state may be entered, and flags a reference that is absent on
every such environment.

\subsection{The concrete property}

Fix a workflow with state set $S$. A concrete execution threads a JSON document \emph{and} a store
of workflow variables (written by \texttt{Assign}); write $\mathcal{D}$ for the set of JSON values.
For a state $s$, let $\llbracket s\rrbracket$ be the set of \emph{environments}---a document paired
with a variable store---with which $s$ is entered over all executions that reach $s$. A path
$p=k_1.k_2.\dots.k_m$ is \emph{present} in a JSON value $d$ if $d$ is an object with key $k_1$ whose
value has $k_2$, and so on to $k_m$; otherwise $p$ is \emph{absent}. Each read site selects one
source---the document, or, for a \texttt{\$}\emph{name} reference, the variable store. The analysis
answers a single question soundly: \emph{is the read path $p$ absent in that selected source on every
environment in $\llbracket s\rrbracket$?} If so, the reference evaluated at $s$ fails on every
execution (a JSONPath path-not-found, an unresolved variable, or a JSONata
\texttt{States.QueryEvaluationError}), and \tool{} reports \code{SC1101}.

\subsection{Abstract domain}

A \emph{document shape} abstracts a set of documents by the field paths that \emph{may} be
present.

\begin{definition}[Shapes]
The set of shapes $\Sigma$ is generated by
\[
  \sigma \;::=\; \top \;\mid\; \Obj\,r, \qquad r : \mathit{Key}\rightharpoonup\Sigma
  \text{ a finite partial map.}
\]
$\top$ is opaque (any field may be present). $\Obj\,r$ constrains an object to the keys in
$\dom(r)$: every top-level key that \emph{may} occur is listed, so a key outside $\dom(r)$ is
present on no represented \emph{object}. The abstraction functions send a scalar to
$\Obj\,\emptyset$ (no fields) and an array to $\top$ (its elements are unconstrained); this is
a choice of the abstraction, not a claim that $\gamma(\Obj\,r)$ contains only objects---indeed
$\gamma(\Obj\,r)$ admits every non-object (Definition~\ref{def:gamma}), so $\Obj\,r$ is a
constraint on the object case alone.
\end{definition}

\begin{definition}[Order and join]
The join $\sqcup$ is
\[
  \top\sqcup\sigma=\sigma\sqcup\top=\top,\qquad
  \Obj\,r_1 \sqcup \Obj\,r_2 = \Obj\,r,\ \
  \dom(r)=\dom(r_1)\cup\dom(r_2),
\]
where $r(k)=r_1(k)\sqcup r_2(k)$ if $k\in\dom(r_1)\cap\dom(r_2)$, and $r(k)=r_i(k)$ if $k$
occurs only in $r_i$. The induced order is $\sigma_1\sqsubseteq\sigma_2 \iff
\sigma_1\sqcup\sigma_2=\sigma_2$; equivalently $\sigma\sqsubseteq\top$ for all $\sigma$, and
$\Obj\,r_1\sqsubseteq\Obj\,r_2$ iff $\dom(r_1)\subseteq\dom(r_2)$ and
$r_1(k)\sqsubseteq r_2(k)$ for all $k\in\dom(r_1)$.
\end{definition}

\noindent Intuitively a shape is \emph{higher} when it admits \emph{more} documents: $\top$
is the top, $\Obj\,\emptyset$ the most precise object shape, and the join is the least upper
bound. The join is exactly the key-union of the implementation (clone the left record, then
for each key of the right record join or insert).

\begin{definition}[Definite absence]
A path $p=k_1.\dots.k_m$ is \emph{definitely absent} in $\sigma$ if descending $\sigma$ by
$k_1,\dots$ reaches some $\Obj\,r$ with $k_i\notin\dom(r)$ \emph{before} encountering $\top$.
Descent stops at $\top$: a path under $\top$ is never definitely absent.
\end{definition}

\subsection{Concretization}

\begin{definition}[Concretization $\gamma$]\label{def:gamma}
$\gamma:\Sigma\to\mathcal{P}(\mathcal{D})$ is
\[
  \gamma(\top)=\mathcal{D},\qquad
  \gamma(\Obj\,r)=\bigl\{\,d\in\mathcal{D}\ \big|\ d \text{ object} \Rightarrow
  \forall k\in\mathit{keys}(d).\ k\in\dom(r)\ \wedge\ d(k)\in\gamma(r(k))\,\bigr\}.
\]
A non-object $d$ (scalar, array) has no top-level keys, so it lies in $\gamma(\Obj\,r)$ for
every $r$.
\end{definition}

\begin{lemma}[Absence is captured]\label{lem:absence}
If $p$ is definitely absent in $\sigma$, then $p$ is absent in every $d\in\gamma(\sigma)$.
\end{lemma}
\begin{proof}
Induction on $|p|$. If $p=k_1.\dots$ and $\sigma=\Obj\,r$ with $k_1\notin\dom(r)$, then by
definition of $\gamma$ no $d\in\gamma(\Obj\,r)$ has top-level key $k_1$, so $p$ is absent in
$d$. Otherwise $k_1\in\dom(r)$ and $k_2.\dots.k_m$ is definitely absent in $r(k_1)$; for any
$d\in\gamma(\Obj\,r)$, either $d$ lacks $k_1$ (so $p$ absent) or $d(k_1)\in\gamma(r(k_1))$,
and by the induction hypothesis $k_2.\dots.k_m$ is absent in $d(k_1)$, so $p$ is absent in
$d$. Definite absence cannot pass through $\top$, so the case $\sigma=\top$ never arises along
the descent.
\end{proof}

\begin{lemma}[$\gamma$ is monotone, $\sqcup$ is sound]\label{lem:join}
$\sigma_1\sqsubseteq\sigma_2\Rightarrow\gamma(\sigma_1)\subseteq\gamma(\sigma_2)$, and
$\gamma(\sigma_1)\cup\gamma(\sigma_2)\subseteq\gamma(\sigma_1\sqcup\sigma_2)$.
\end{lemma}
\begin{proof}
Monotonicity is by induction on shape structure: if $\sigma_2=\top$ it is immediate; for
$\Obj\,r_1\sqsubseteq\Obj\,r_2$ we have $\dom(r_1)\subseteq\dom(r_2)$ and
$r_1(k)\sqsubseteq r_2(k)$, so any $d$ admitted by $\Obj\,r_1$ has $\mathit{keys}(d)\subseteq
\dom(r_1)\subseteq\dom(r_2)$ and $d(k)\in\gamma(r_1(k))\subseteq\gamma(r_2(k))$ by the
induction hypothesis. Soundness of $\sqcup$ follows since
$\sigma_i\sqsubseteq\sigma_1\sqcup\sigma_2$ and $\gamma$ is monotone.
\end{proof}

\subsection{Environments: document and variables}
\label{sec:tr:env}

Current \asl{} threads not one but \emph{two} pieces of state: the JSON document and a store of
workflow \emph{variables}, written by \texttt{Assign} and read as \texttt{\$}\emph{name} (JSONPath
mode) or as \texttt{\$}\emph{name} inside a \texttt{\{\%\ \%\}} expression (JSONata mode). Because a
field can now travel through a variable instead of the document, provenance must track the variable
store to stay sound: a read the analysis cannot see satisfied would otherwise be flagged as a
missing document field. We abstract the store by a \emph{single} shape---the record whose keys are
the variable names that may be bound, each mapped to its own shape---drawn from the \emph{same}
lattice $\Sigma$. An abstract \emph{environment} is then a pair
\[
  \Env \;=\; \Sigma \times \Sigma, \qquad e = (\sigma_{\mathrm d},\,\sigma_{\mathrm v}),
\]
with document component $\sigma_{\mathrm d}$ and variable-store component $\sigma_{\mathrm v}$.
Order, join, and concretization are componentwise:
$(\sigma_{\mathrm d},\sigma_{\mathrm v})\sqsubseteq(\sigma_{\mathrm d}',\sigma_{\mathrm v}')$ iff
$\sigma_{\mathrm d}\sqsubseteq\sigma_{\mathrm d}'$ and
$\sigma_{\mathrm v}\sqsubseteq\sigma_{\mathrm v}'$;
$(\sigma_{\mathrm d},\sigma_{\mathrm v})\sqcup(\sigma_{\mathrm d}',\sigma_{\mathrm v}')=
(\sigma_{\mathrm d}\sqcup\sigma_{\mathrm d}',\ \sigma_{\mathrm v}\sqcup\sigma_{\mathrm v}')$; and
\[
  \gamma_{\Env}(\sigma_{\mathrm d},\sigma_{\mathrm v})=
  \bigl\{\,(d,\nu)\ \big|\ d\in\gamma(\sigma_{\mathrm d}),\ \nu\in\gamma(\sigma_{\mathrm v})\,\bigr\},
\]
a concrete environment $(d,\nu)$ pairing the document $d$ with the variable store $\nu$ (itself a
JSON object mapping names to values). The empty store $\Obj\,\emptyset$ (no variable bound) seeds
$\sigma_{\mathrm v}$. Because $\Env$ is a finite product of the finite-height lattice $\Sigma$, it
is again finite-height, and Lemmas~\ref{lem:absence} and~\ref{lem:join} lift componentwise. All
transfers below act on environments; one that touches only the document leaves $\sigma_{\mathrm v}$
unchanged, and \texttt{Assign} touches only $\sigma_{\mathrm v}$.

\subsection{Transfer functions}

Each stage is a monotone function on shapes. We write $\narrow(\sigma,\pi)$ for re-rooting
along a path expression $\pi$ (used by \texttt{InputPath}/\texttt{OutputPath}),
$\mathit{ctor}(v)$ for the record built from a constructor value $v$ (used by
\texttt{Parameters}/\texttt{ItemSelector}/\texttt{ResultSelector}), and
$\place(\sigma,\rho,\sigma_r)$ for merging a result shape $\sigma_r$ at a
\texttt{ResultPath} $\rho$.

\paragraph{Re-rooting.}
Write $\varepsilon$ for the empty (absent) path and $k\!::\!\mathit{rest}$ for the split of a
path into its first key and remainder. Then $\narrow(\sigma,\varepsilon)=\sigma$ (no path) and
$\narrow(\sigma,\$)=\sigma$ (the bare root); $\narrow(\sigma,\texttt{null})=\Obj\,\emptyset$,
sound because \asl{} maps an \texttt{InputPath}/\texttt{OutputPath} of \texttt{null} to the
empty object $\{\}\in\gamma(\Obj\,\emptyset)$. For a precise member path
$\pi=k_1.\dots.k_n$---written either in dotted form or quoted-bracket member form---$\narrow$
descends via $\mathit{nav}$: $\mathit{nav}(\top,\_)=\top$ and
$\mathit{nav}(\Obj\,r,k\!::\!\mathit{rest})=\mathit{nav}(r(k),\mathit{rest})$ if $k\in\dom(r)$,
and $=\top$ otherwise (a missing intermediate key yields $\top$, over-approximating). Any path
outside this member-path fragment (wildcards, filters, functions, array indexes, the \texttt{\$\$}
context) yields $\top$; a non-string, non-\texttt{null} path value is treated as the identity
($\narrow(\sigma,\cdot)=\sigma$), a case \asl{} deploy-time validation rejects and which is
therefore unreachable in practice.

\paragraph{Constructors and literals.}
Two builders recurse over a JSON value $v$. The \emph{constructor} $\mathit{ctor}(v)$ models a
payload template (\texttt{Parameters}/\texttt{ItemSelector}/\texttt{ResultSelector}): for an
object it builds $\Obj\,r$ where a key ``$k$\texttt{.\$}'' contributes $r(k)=\top$ (it
\emph{reads} an opaque value into field $k$), a static key $k$ contributes
$r(k)=\mathit{ctor}(v(k))$, a scalar leaf is $\Obj\,\emptyset$, and an array is $\top$. The
\emph{literal} builder $\mathit{lit}(v)$ models verbatim data (a \texttt{Pass} \texttt{Result}),
which \asl{} copies uninterpreted: it agrees with $\mathit{ctor}$ \emph{except} that it does no
\texttt{.\$} stripping---a key literally named ``$k$\texttt{.\$}'' stays a key of that exact
name. Distinguishing the two is necessary for soundness: for
\texttt{Result}$=\{$``\texttt{a.\$}''$:v\}$ the concrete document has the top-level key
``\texttt{a.\$}'', which is \emph{not} in $\dom(\mathit{ctor})=\{a\}$, so $\mathit{ctor}$ would
under-approximate the key set.

\paragraph{Result and merge.}
The result shape $\sigma_{\mathit{res}}$ of a state, built from its effective input $\eff$, is:
$\top$ for an opaque \texttt{Map}/\texttt{Parallel} result and for a \texttt{Task} result with no
modeled result shape; $\mathit{ctor}(\textsf{ResultSelector})$ when a \texttt{ResultSelector} pins
it; a closed record when the optional result-shape model supplies a declared task-output contract or
a fixed AWS service-result envelope; and for a \texttt{Pass}, $\mathit{lit}(\textsf{Result})$ if it
carries a literal \texttt{Result}, else $\mathit{ctor}(\textsf{Parameters})$ if it carries only
\texttt{Parameters}, else $\eff$ itself (a bare \texttt{Pass} forwards its effective input
unchanged). Merging is
\[
  \place(\sigma,\rho,\sigma_r)=
  \begin{cases}
    \sigma_r & \rho=\textsf{Default}\ \text{or}\ \rho=\$ \quad(\text{result replaces the document}),\\
    \sigma & \rho=\textsf{Discard}\ (\texttt{null})\quad(\text{result dropped, input kept}),\\
    \place_{\mathit{in}}(\sigma,\pi,\sigma_r) & \rho=\textsf{Path}(\pi),
  \end{cases}
\]
where $\place_{\mathit{in}}$ descends $\pi$, cloning each record and \emph{preserving
siblings}, and writes $\sigma_r$ at the leaf; a $\top$ parent stays $\top$, and an
unparseable $\pi$ yields $\top$. (This is why a field can survive a task that never mentions
it: a sub-path merge leaves the rest of the document intact.)

\paragraph{Variable assignment.}
An \texttt{Assign} block updates the variable store, \emph{overwriting} any earlier binding of the
same name (bindings from distinct control-flow paths are joined only where they merge, at the
fixpoint). Write $\bind(\sigma_{\mathrm v},k,\sigma)$ for the store that agrees with
$\sigma_{\mathrm v}$ except that key $k$ maps to $\sigma$ (identity when $\sigma_{\mathrm v}=\top$).
In JSONPath mode an entry ``$k$\texttt{.\$}''$:\pi$ binds
$k\mapsto\narrow(\sigma_{\mathrm{src}},\pi)$ and a static entry $k:v$ binds $k\mapsto\mathit{ctor}(v)$,
where the source shape $\sigma_{\mathrm{src}}$ is the value \asl{} reads an \texttt{Assign} against
for that state kind---the raw result ($\top$, unless result-shape modeling supplies a task result
record) for \texttt{Task}, $\top$ for \texttt{Map}/\texttt{Parallel}, the \texttt{Pass} result
shape, or the effective input $\eff$ for \texttt{Choice}/\texttt{Wait}. In
JSONata mode an entry $k:v$ binds $k\mapsto\jval(v)$ (below). The whole block $\asgn(\sigma_{\mathrm
v},\textsf{Assign})$ is the composition of these single-key binds.

\paragraph{JSONata expressions.}
A JSONata value---a string ``\texttt{\{\%}\,$e$\,\texttt{\%\}}'', or a JSON object/array containing
such strings---denotes a shape $\jval$ against the current sources: the entry document, the state's
raw result, an opaque error object, and the variable store, each itself a shape. On the
\emph{modeled fragment} a \emph{direct reference} \texttt{\$states.input}$.\pi$,
\texttt{\$states.result}$.\pi$, \texttt{\$states.errorOutput}$.\pi$, or \texttt{\$}\emph{name}$.\pi$
navigates ($\mathit{nav}$) the corresponding source shape at $\pi$; an \emph{object literal}
$\{\,\text{'}k_1\text{'}\!:e_1,\dots\}$ builds $\Obj\{k_i\mapsto\jval(e_i)\}$; a surrounding JSON
object/array recurses into its members. Every other expression---functions such as
\texttt{\$exists}, filters, arithmetic, string operators, and \texttt{\$states.context}---yields
$\top$. A modeled read thus resolves against exactly one source shape and an unmodeled one is opaque:
the JSONata analogue of the simple-path fragment. This is what lets a field move through a variable
in JSONata mode (\texttt{Assign}$\{k:\text{\{\%}\,\texttt{\$states.input}.\pi\,\text{\%\}}\}$ then a
later \texttt{\$}$k$ read) without a spurious finding, while still flagging a read of a name no path
ever assigns.

\paragraph{State transfer and error edges.}
For a state with entry environment $e=(\sigma_{\mathrm d},\sigma_{\mathrm v})$, the output
environment $\mathit{out}(e)=(\sigma_{\mathrm d}',\sigma_{\mathrm v}')$ updates both components. The
variable store becomes $\sigma_{\mathrm v}'=\asgn(\sigma_{\mathrm v},\textsf{Assign})$ (identity when
the state has no \texttt{Assign}). The document becomes
\[
  \sigma_{\mathrm d}'=
  \begin{cases}
    \jout(e) & \text{JSONata mode (query state),}\\
    \narrow\!\bigl(\place(\sigma_{\mathrm d},\rho,\sigma_{\mathit{res}}),\ \textsf{OutputPath}\bigr)
      & \texttt{Task}/\texttt{Pass}/\texttt{Map}/\texttt{Parallel},\\
    \narrow(\eff,\ \textsf{OutputPath}) & \text{otherwise,}
  \end{cases}
\]
with $\eff=\narrow(\sigma_{\mathrm d},\textsf{InputPath})$ the \texttt{InputPath}-selected input from
which the result shape $\sigma_{\mathit{res}}$ is built, and $\rho$ the state's \texttt{ResultPath}.
The JSONPath result is merged into the \emph{raw} entry document $\sigma_{\mathrm d}$, \emph{not}
$\eff$---exactly as \asl{} writes \texttt{ResultPath} into the pre-\texttt{InputPath} document, so a
field the \texttt{InputPath} selection dropped still survives. In JSONata mode the successor document
is the shape of the state's \texttt{Output}, $\jout(e)=\jval(\textsf{Output})$---an object literal
closes the successor shape precisely, and a query \texttt{Task} with no \texttt{Output} yields its
modeled raw-result shape when result-shape modeling is enabled, otherwise $\top$ (as do
\texttt{Map}/\texttt{Parallel}); \asl{}'s JSONata mode has no \texttt{InputPath}/\texttt{OutputPath}
to re-root. A \texttt{Catch} edge from entry $e$ propagates
$\mathit{catch}(e,\rho)=\bigl(\place(\sigma_{\mathrm d},\rho,\top),\ \asgn(\sigma_{\mathrm
v},\textsf{Catch.Assign})\bigr)$: the error object is opaque ($\top$), merged at the catch's
\texttt{ResultPath}, and the catcher's own \texttt{Assign} (if any) updates the store. A
\texttt{Retry} re-executes the same state on the same entry environment, so it leaves entry shapes
unchanged and the analysis ignores it.

\begin{lemma}[Local soundness]\label{lem:local}
Writing a subscript $\mathrm{c}$ for the concrete counterpart of each stage, every transfer is
sound on its own arguments:
\begin{itemize}
\item[(re-root)] $d\in\gamma(\sigma)\Rightarrow \narrow_{\mathrm{c}}(d,\pi)\in\gamma(\narrow(\sigma,\pi))$;
\item[(build)] the concrete payload built over a document $d$ from a template $v$ lies in
$\gamma(\mathit{ctor}(v))$, and from a literal $v$ in $\gamma(\mathit{lit}(v))$;
\item[(merge)] $d\in\gamma(\sigma)\wedge d_r\in\gamma(\sigma_r)\Rightarrow
\place_{\mathrm{c}}(d,\rho,d_r)\in\gamma(\place(\sigma,\rho,\sigma_r))$;
\item[(jsonata)] if each source value lies in the $\gamma$ of its source shape, the concrete value
of a JSONata expression $e$ lies in $\gamma(\jval(e))$;
\item[(assign)] $\nu\in\gamma(\sigma_{\mathrm v})\Rightarrow
\asgn_{\mathrm c}(\nu,\textsf{Assign})\in\gamma(\asgn(\sigma_{\mathrm v},\textsf{Assign}))$;
\item[(state/catch)] $(d,\nu)\in\gamma_{\Env}(e)\Rightarrow
\mathit{out}_{\mathrm{c}}(d,\nu)\in\gamma_{\Env}(\mathit{out}(e))$ and
$\mathit{catch}_{\mathrm{c}}(d,\nu)\in\gamma_{\Env}(\mathit{catch}(e,\rho))$.
\end{itemize}
\end{lemma}
\begin{proof}
By cases. \emph{Re-rooting:} if $\sigma=\top$ or $\pi$ is non-simple, $\narrow(\sigma,\pi)=\top$
and $\gamma(\top)=\mathcal{D}$; for a simple path, $\mathit{nav}$ keeps $d(k)\in\gamma(r(k))$ by
Definition~\ref{def:gamma}, and a missing intermediate key yields $\top$ (vacuously sound);
$\narrow(\sigma,\texttt{null})=\Obj\,\emptyset$ is sound because the selected concrete value is
the empty object $\{\}$, which has no keys and hence lies in $\gamma(\Obj\,\emptyset)$.
\emph{Constructor:} the concrete payload object has top-level keys exactly the template's keys
$\subseteq\dom(\mathit{ctor}(v))$; a \texttt{.\$} key binds a runtime-read value in
$\gamma(\top)$, and a static key's value lies in $\gamma(\mathit{ctor}(v(k)))$ by induction.
\emph{Literal:} the concrete value is $v$ copied verbatim, whose keys are exactly
$\dom(\mathit{lit}(v))$---including any literal ``\texttt{.\$}'' key---and each child lies in the
corresponding $\gamma$ by induction. \emph{Result:} an opaque result lies in $\gamma(\top)$; a
pinned \texttt{ResultSelector} reduces to the constructor case; a bare \texttt{Pass} forwards
$\eff$, and $d\in\gamma(\eff)$ by hypothesis. \emph{Merge:} $\rho\in\{\textsf{Default},\$\}$
replaces, and the concrete output is the result, which lies in $\gamma(\sigma_r)$;
$\rho=\textsf{Discard}$ keeps the input; for $\rho=\textsf{Path}(\pi)$, $\place_{\mathit{in}}$
writes $\sigma_r$ at $\pi$ and preserves siblings, matching the concrete merge, and a $\top$
parent (or unparseable $\pi$) stays $\top$ (sound). \emph{JSONata:} a modeled direct reference
navigates one source shape and is sound by the re-root case ($\mathit{nav}$ is $\narrow$'s
descent); an object literal is the constructor case over $\jval$-typed children; every unmodeled
expression yields $\top$, whose $\gamma$ admits the concrete value. \emph{Assign:} each $\bind$
replaces one store key by a shape that---by the (build), (jsonata), or (re-root) case---
over-approximates the concrete bound value while leaving the other keys' admission unchanged, so
the updated store lies in $\gamma$ of the updated store shape; composing the entries yields
$\asgn$. \emph{State transfer} pairs the document update (the JSONPath composition above, or
$\jout$ in JSONata mode) with the variable update $\asgn$; since $\gamma_{\Env}$ is the
componentwise conjunction, soundness of each component gives soundness of the pair. \emph{Catch}
places $\top$ (the error object) at $\rho$ in the document---sound since
$\gamma(\top)=\mathcal{D}$---and applies $\asgn$ for the catcher's \texttt{Assign}.
\end{proof}

\subsection{The fixpoint}

The analysis maintains a \emph{partial} map $\Phi:S\rightharpoonup\Env$, seeding only the
start state with $e_0=(\sigma_0,\Obj\,\emptyset)$ (document seed as below, empty variable store);
a state acquires an environment the first time an edge reaches it, and a state never reached stays
outside $\dom(\Phi)$ and is read as $(\top,\top)$ at check time (an unreached state constrains
nothing, and reporting against $\top$ never fires). One \emph{round}
updates, for every $s\in\dom(\Phi)$ with current entry $\Phi(s)$: for each normal successor
$t$, $\Phi(t)\leftarrow\Phi(t)\sqcup\mathit{out}(\Phi(s))$ (inserting $\mathit{out}(\Phi(s))$
when $t\notin\dom(\Phi)$); for each \texttt{Catch} target $u$ with result path $\rho$,
$\Phi(u)\leftarrow\Phi(u)\sqcup\mathit{catch}(\Phi(s),\rho)$. Updates are in place
(Gauss--Seidel) in an unspecified state order; the least fixpoint, once reached, is independent
of that order. Rounds repeat until a round changes nothing.

\begin{lemma}[Widened convergence]\label{lem:term}
The forward iteration with depth-$k$ widening is monotone and reaches a post-fixpoint on every
workflow topology. The implementation uses $k=12$ and records the safety bound
$N=|S|\cdot(\kappa_0+2)\cdot(k+1)+2$ rounds, where $\kappa_0$ is the implementation's key count.
If the safety bound were exceeded because of an implementation bug, all known entries would be
lifted to $(\top,\top)$ before checking, preserving soundness by losing precision rather than
reporting from an under-approximation.
\end{lemma}
\begin{proof}
Every normal or catch-edge update is a join, so each $\Phi(s)$ is non-decreasing
(Lemma~\ref{lem:join}). Let $\kappa$ be the workflow's finite \emph{key universe}: keys occurring
in constructors (\texttt{Parameters}, \texttt{Result}, \texttt{ResultSelector},
\texttt{ItemSelector}), in \texttt{Assign}/\texttt{Arguments}/\texttt{Output} templates or JSONata
object literals, in declared input or output schemas, and as segments of \texttt{ResultPath} and
catch \texttt{ResultPath} fields. Define $W_k$ over shapes by replacing every record nested deeper
than $k$ by $\top$ and recursing otherwise. $W_k$ is extensive and monotone: for every shape
$\sigma$, $\gamma(\sigma)\subseteq\gamma(W_k(\sigma))$, and if $\sigma_1\sqsubseteq\sigma_2$ then
$W_k(\sigma_1)\sqsubseteq W_k(\sigma_2)$. The implementation applies $W_k$ to every propagated
environment and again after each join.

After widening, every environment component ranges over records of depth at most $k$ over the
finite key universe, plus $\top$; the product environment lattice therefore has finite height. A
monotone in-place iteration over a finite-height lattice stabilises at a post-fixpoint independent
of state order. Cyclic re-embedding can still push a growing tail through a loop, but once the tail
passes depth $k$ it becomes $\top$, so the chain stops growing. The emergency lift to $(\top,\top)$
is also an over-approximation: subsequent propagation can only join smaller transfer results into
$\top$ entries, so no definite-absence report can arise from the lifted region.
\end{proof}

\begin{remark}[Empirical headroom]
The widened bound is deliberately conservative. On the current corpus---193 AWS workflows, 66 CNCF
workflows, and two DSL workflows, including nested machines and both native and typed seeding---all
790 data-flow fixpoints converge, with a maximum of 7 rounds and a worst observed bound utilization
of 2/54 (3.7\%). The synthetic scale sweep remains near-linear; acyclic chains converge in one
forward pass by construction.
\end{remark}

\subsection{Reference checking}

At the widened fixpoint (Lemma~\ref{lem:term}), the analysis checks,
at each state $s$ with entry environment $\Phi(s)=(\sigma_{\mathrm d},\sigma_{\mathrm v})$ and
effective document $\eff=\narrow(\sigma_{\mathrm d},\textsf{InputPath})$, the reads the runtime
evaluates and aborts on when absent. These fall into three classes.

\emph{(i) JSONPath payload reads.} The \texttt{.\$} keys (at any depth) of \texttt{Parameters} and
\texttt{ItemSelector} and a \texttt{Map}'s \texttt{ItemsPath}, resolved against $\eff$; and the
\texttt{.\$} keys of a JSONPath \texttt{Assign}, resolved against that state's assign source. A
reference \texttt{\$}\emph{name}$.\pi$ that names a workflow \emph{variable} (a leading \texttt{\$}
not followed by ``\texttt{.}'', hence not a document path \texttt{\$.\dots}) is resolved against the
store $\sigma_{\mathrm v}$ rather than $\eff$---so a field that travelled through a variable is seen
where it is read, and only a variable no path assigns is flagged.

\emph{(ii) Modeled JSONata reads.} In a JSONata state, a field-value string whose \emph{entire}
body is a direct reference (\texttt{\$states.input}$.\pi$, \texttt{\$states.result}$.\pi$,
\texttt{\$states.errorOutput}$.\pi$, or \texttt{\$}\emph{name}$.\pi$)---in \texttt{Arguments},
\texttt{Output}, \texttt{Assign}, or a \texttt{Choice} \texttt{Condition}---resolved against its
source shape (effective document, raw result, error object, or store). Such an expression returns
\texttt{undefined} on a missing path, which \asl{} raises as
\texttt{States.QueryEvaluationError}~\cite{awsasl} (``failure to return a result'')---a definite
runtime fault, exactly the \code{SC1101} contract. A reference nested inside a JSONata \emph{object
constructor} $\{\,\text{'}k\text{'}\!:e\}$ is deliberately \emph{not} in this class: JSONata omits a
key whose value is undefined, so the constructor yields a value and does not fail; such a reference
shapes the successor document (\S\ref{sec:tr:df}) but is never itself flagged.

\emph{(iii) Excluded (opaque) reads.} Any string not beginning with \texttt{\$}, the context object
\texttt{\$\$}/\texttt{\$states.context}, \texttt{States.*} intrinsics, the bare root \texttt{\$}, any
JSONPath outside the simple-path fragment, and any JSONata expression outside the modeled
direct-reference/object fragment (functions such as \texttt{\$exists}, filters, arithmetic). All
resolve to $\top$ and never fire.

For a checked reference to a path $p$ against its selected source shape $\sigma$ (document or
variable store), \tool{} reports \code{SC1101} iff $p$ is definitely absent in $\sigma$ (the lookup
returns \textsf{Missing}; a path into a $\top$ or otherwise unconstrained region returns
\textsf{Maybe} and is never flagged). A JSONPath \texttt{Choice} operand is handled as a
\emph{warning}, not an error; see Section~\ref{sec:tr:deadguard}.

\begin{assumption}[Runtime resolution on non-objects]\label{asm:jsonpath}
The \asl{} JSONPath engine resolves a dotted reference $k_1.\dots.k_m$ by successive object-key
lookup: applying a key to a non-object (scalar or array) fails with a path-not-found error, and
a numeric segment denotes a property name rather than an array index~\cite{awsasl,jsonpath}.
Under this assumption the concrete \emph{present} relation of Section~\ref{sec:tr:df}
(object-key descent) coincides with runtime resolution, and $\gamma$'s admission of arbitrary
non-objects under $\Obj\,r$ is sound: a reference descending into a non-object simply fails,
exactly as ``absent'' predicts.
\end{assumption}

\begin{theorem}[No false positives]\label{thm:soundness}
If \tool{} reports \code{SC1101} for a path $p$ at a state $s$, then in every concrete execution
that reaches $s$, $p$ is absent from the document \emph{or variable store} that the read site
selects, so the \asl{} runtime fails the reference---a JSONPath path-not-found, an unresolved
variable, or a JSONata \texttt{States.QueryEvaluationError}. Equivalently, the analysis never flags
a reference some execution could satisfy.
\end{theorem}
\begin{proof}
First, the fixpoint over-approximates the reachable environments: for every $s$ and every reachable
$(d,\nu)\in\llbracket s\rrbracket$ (document $d$ paired with variable store $\nu$),
$(d,\nu)\in\gamma_{\Env}(\Phi(s))$. Induction on the length of the execution prefix reaching $s$.
\emph{Base:} the seed $e_0=(\sigma_0,\Obj\,\emptyset)$ satisfies it---the document part as before
(native tier $\sigma_0=\top$ and $\gamma(\top)=\mathcal{D}$; typed tier $\sigma_0=\Obj\,r_0$ over the
declared input fields, under the closed-world reading of the contract), and the initial store is
empty, so $\nu_0=\{\}\in\gamma(\Obj\,\emptyset)$. Closed task-result records are sound relative to
the declared output contract or fixed AWS service envelope they model. \emph{Step:} suppose $s$ is
reached from $s'$. On a
normal transition the concrete environment is $\mathit{out}_{\mathrm c}(d',\nu')$ for some
$(d',\nu')\in\llbracket s'\rrbracket$; by the induction hypothesis
$(d',\nu')\in\gamma_{\Env}(\Phi(s'))$, by Lemma~\ref{lem:local} (state/catch, which composes the
document transfer with the \texttt{Assign} store update)
$\mathit{out}_{\mathrm c}(d',\nu')\in\gamma_{\Env}(\mathit{out}(\Phi(s')))$, and since the round
joined $\mathit{out}(\Phi(s'))$ into $\Phi(s)$, by the componentwise lift of Lemma~\ref{lem:join}
the concrete environment lies in $\gamma_{\Env}(\Phi(s))$. A \texttt{Catch} transition is identical
using $\mathit{catch}$, which places $\top$ in the document and applies the catcher's \texttt{Assign}.
The \texttt{InputPath} narrowing at $s$ is sound (Lemma~\ref{lem:local}), so effective documents lie
in $\gamma(\eff)$ and stores in $\gamma(\sigma_{\mathrm v})$.

Now suppose \code{SC1101} fires for $p$ at $s$, against its selected source shape $\sigma$---the
effective document $\eff$ for a document read, or the store $\sigma_{\mathrm v}$ for a variable read.
The check fires only when $p$ is definitely absent in $\sigma$, so by Lemma~\ref{lem:absence} no
value in $\gamma(\sigma)$ contains $p$; since every concrete source at $s$ lies in $\gamma(\sigma)$,
$p$ is absent on all of them. A JSONPath payload read then fails by Assumption~\ref{asm:jsonpath}
(path-not-found); a JSONPath variable read has no binding to resolve; and a modeled JSONata direct
reference evaluates to \texttt{undefined}, which \asl{} raises as \texttt{States.QueryEvaluationError}.
In every case the runtime aborts the reference. A finding is emitted only when the machine converged
(Lemma~\ref{lem:term}), i.e.\ at a genuine fixpoint, so the conclusion holds whenever \code{SC1101}
is reported.
\end{proof}

\begin{corollary}[Native unconditionality]
On the native tier ($\sigma_0=\top$), Theorem~\ref{thm:soundness} holds for an arbitrary
execution input. On the typed tier it holds relative to the closed-world reading of the
declared input schema; when result-shape modeling is enabled it also holds relative to the
closed-world reading of declared task-output schemas and fixed AWS service-result envelopes. Where
inputs or results are open, the native opaque shape is the sound fallback. (The typed seed also
admits a second closed-world source: when a workflow declares no input schema but its start task
declares an input contract, the start seed is that contract's field record.)
\end{corollary}

\subsection{Nested machines}
\label{sec:tr:nested}

A \texttt{Map}'s iterator and a \texttt{Parallel}'s branches are themselves state machines, each
analysed by the same procedure under a seed derived from the enclosing state's entry.

\begin{remark}[Nested-machine seeds]\label{rmk:nested}
The nested analysis runs under seeds derived from the enclosing machine's widened fixpoint; if the
safety guard ever lifted an enclosing entry to $\top$, the nested seed is correspondingly opaque and
conservative. Each \texttt{Map}
iteration is seeded with $\mathit{ctor}(\textsf{ItemSelector})$ when an \texttt{ItemSelector} is
declared and $\top$ otherwise; each \texttt{Parallel} branch is seeded with
$\mathit{ctor}(\textsf{Parameters})$ when the state declares \texttt{Parameters} and with the
effective input $\eff=\narrow(\Phi(s),\textsf{InputPath})$ otherwise---exactly the document
\asl{} hands each branch. Each seed over-approximates the corresponding per-iteration/per-branch
document (Lemma~\ref{lem:local}), so Theorem~\ref{thm:soundness} extends unconditionally to
findings emitted inside a nested machine.
\end{remark}

\subsection{The dead-guard refinement (\code{SC1110})}
\label{sec:tr:deadguard}

The same shapes support a path-sensitive \emph{warning}. Unlike a \texttt{.\$} payload
reference, a \texttt{Choice} comparator is a \emph{guard} evaluated in rule order: the runtime
tries each rule until one matches and follows its branch. A value comparator on an unresolvable
operand does not silently return false---\asl{} raises a \texttt{States.Runtime}
error~\cite{awsasl}, which is precisely why \texttt{IsPresent} exists. We nonetheless report only
a \emph{warning}, not an \code{SC1101} error, for two reasons: the rule may never be reached (an
earlier rule can match first and divert control), so we cannot promise a runtime abort; and the
finding's real content is that a branch whose guard \emph{requires} a definitely absent field
can never be taken---dead logic. \code{SC1110} fires when a \texttt{Choice} rule's condition
requires presence of a field definitely absent in the state's effective (post-\texttt{InputPath})
shape. Presence is required by value comparators, \texttt{IsNull}, the type predicates, and
\texttt{IsPresent:true}; it is \emph{not} required by \texttt{IsPresent:false}. The check is
conservative on Boolean structure: a conjunction is dead if any conjunct is dead, a disjunction
only if \emph{all} disjuncts are dead, and a negation is never flagged. Soundness of the
underlying ``definitely absent'' judgment is exactly Lemma~\ref{lem:absence}.

\subsection{Coverage of the modeled fragment}
\label{sec:tr:coverage}

The analysis resolves precisely exactly the reads inside the member-path fragment ($\$$, $\$.a$,
$\$.a.b$, $\$[\text{'}a\text{'}]$); every operand outside it lifts to $\top$ and can never produce a finding. To quantify
how much of real ASL this covers, the \texttt{stepcheck path-coverage} command classifies every
field-read a workflow evaluates---the \texttt{.\$} operands of
\texttt{Parameters}/\texttt{ItemSelector}/\texttt{Assign} and a \texttt{Map}'s \texttt{ItemsPath},
i.e.\ the read sites the \code{SC1101} check consumes---into the dotted/quoted-member fragment, the whole-document
$\$$, complex JSONPath (wildcard/filter/function/array-index), the \texttt{\$\$} context object,
\texttt{\$}\emph{name} workflow-variable reads, and \texttt{States.*} intrinsics.
Table~\ref{tab:tr:coverage} reports the split over the AWS corpus, the CNCF examples, and the
independent repositories of Section~\ref{sec:tr:external}. JSONata \texttt{\{\%\ \%\}} reads form a
separate modeled fragment---direct references and object literals (\S\ref{sec:tr:df})---and are not
counted in this JSONPath surface.

\begin{table}[t]
\centering\small
\setlength{\tabcolsep}{5pt}
\caption{Field-read coverage by the modeled JSONPath fragment (\texttt{stepcheck path-coverage}).
``Doc.\ reads'' are the genuine document-field references (member $+$ whole-$\$$ $+$ complex),
excluding the \texttt{\$\$} context and \texttt{States.*} operands, which are not document
references and cannot miss a field. ``Precise'' is dotted or quoted-member access $+$ whole-$\$$.}
\label{tab:tr:coverage}
\begin{tabular}{@{}lrrrrrrrr@{}}
\toprule
Corpus & Reads & Dotted & Whole-$\$$ & Complex & \texttt{\$\$} & \texttt{\$}var & \texttt{States.*} & Precise / doc.\ reads \\
\midrule
AWS (193)        & 1168 & 596 & 146 & 61 & 225 & 0 & 140 & 742/803 = \textbf{92.4\%} \\
CNCF (66)        &    2 &   0 &   0 &  0 &   0 & 0 &   2 & --- (runtime expr.) \\
Independent (95) &  420 & 296 &  57 & 17 &  20 & 0 &  30 & 353/370 = \textbf{95.4\%} \\
\midrule
Combined         & 1590 & 892 & 203 & 78 & 245 & 0 & 172 & 1095/1173 = \textbf{93.4\%} \\
\bottomrule
\end{tabular}
\end{table}

Of the genuine document-field references, $92$--$95\%$ fall inside the modeled member fragment;
only $6.6\%$ (combined) use wildcard, filter, array-index, or other nonmember syntax that the
analysis soundly abstracts to $\top$. No workflow in these JSONPath corpora issues a \texttt{\$}\emph{name} variable read (the
\texttt{\$}var column is zero); the variable and JSONata-fragment provenance is exercised instead by
the regression suite and the JSONata-mode workflows in the corpus. The CNCF corpus exercises almost no \texttt{.\$} reads because Serverless Workflow uses
runtime expressions rather than JSONPath provenance, consistent with the ASL-specific data-flow
checks staying silent there.

\subsection{External validity: independent repositories}
\label{sec:tr:external}

Beyond the three aws-samples collections that make up the main corpus, we ran \tool{} unmodified on
\textbf{95} ASL definitions mined (\texttt{eval/mine\_wild.js}) from \textbf{16} public repositories
outside them---50 domain workflows across 13 repositories, the remainder validator/library
fixtures. Two findings from the sound and conservative native tiers illustrate that the analyses are
not overfit to AWS's samples.

\paragraph{A sound \code{SC1101} in the wild.} In a credit-card \emph{transaction-reversal} workflow,
a \texttt{Map} state \texttt{ApplyReversalEntries} iterates \texttt{ItemsPath: \$.ledgerEntries}. The
start state \texttt{SetReversalDefaults} is a \texttt{Pass} whose \texttt{Parameters} reconstruct the
document as a closed record of seven fields (\texttt{transactionId}, \texttt{amount},
\texttt{currency}, \texttt{accountId}, \texttt{reason}, \texttt{reversalId}, \texttt{reversalStatus});
\texttt{ledgerEntries} is not among them, and the two intervening states merge only their own results
under \texttt{\$.validation} and \texttt{\$.screening}. By the transfer functions of
Section~\ref{sec:tr:df}, the shape at \texttt{ApplyReversalEntries} is a closed record with
$\texttt{ledgerEntries}\notin\dom$, so $\mathrm{present}(\$.ledgerEntries)$ is definitely false and
\code{SC1101} fires soundly---at run time the \texttt{Map} fails to resolve \texttt{ItemsPath}.
Because \texttt{SetReversalDefaults} \emph{reconstructs} the document rather than augmenting it, even
a caller that supplies \texttt{ledgerEntries} has it dropped; an adjacent-task field-subset check,
which never models the \texttt{Parameters} reconstruction, cannot see this.

\paragraph{A conservative \code{SC5001} in the wild.} In an IoT firmware-upgrade workflow, a
\texttt{Parallel} state runs three branches that each write the DynamoDB table \texttt{FWUTaskTokens}
(\texttt{PutItem}); the branches receive document copies, so the shared external table is the only
race, and the concurrency pass reports the provable last-writer-wins collision.

\paragraph{Declared-tier evidence.} The declared tier is silent on the wild corpus only because no
workflow ships a compensator. Supplying the missing facts recovers it: for an independent
payment-settlement pipeline, a sidecar declaring the payment authorization and the settlement as
persistent effects (both are durable money movements) promotes \tool{}'s finding from an inferred
warning to a declared-tier \emph{error}---\code{SC4001} on the uncompensated \texttt{FlagFraudulentPayment}
and \code{SC4010} on \texttt{AuthorizePayment}/\texttt{SettlePayment}, whose \texttt{Catch} paths only
reach a \texttt{Fail} state, so a settlement failure leaves an authorized hold that is never voided.
The sidecar and reproduction command ship in the artifact (\texttt{corpus/wild-annot/}).

\section{The workflow type discipline}
\label{sec:tr:types}

\subsection{Signatures, judgment, and validity}

Let $\Gamma$ collect the declared task signatures. Each task $T$ has a \emph{contract}
$A\rightarrow B$ (it consumes the field-record $A$ and produces $B$) and a
\emph{typestate}~\cite{stromYemini,delineFahndrich} transition
$s_{\mathrm{in}}^{T}\!\to\!s_{\mathrm{out}}^{T}$ over business states $S_{\mathrm B}$;
we bundle the two into one signature $T:(A,s_{\mathrm{in}}^T)\to(B,s_{\mathrm{out}}^T)$. The
well-typedness judgment $\Gamma\vdash W:(A,s)\Rightarrow(B,s')$ carries a \emph{pair} type---a
field-record and a business state---and is generated by:
\begin{gather*}
\frac{(T:(A,s_{\mathrm{in}})\to(B,s_{\mathrm{out}}))\in\Gamma}
     {\Gamma\vdash T:(A,s_{\mathrm{in}})\Rightarrow(B,s_{\mathrm{out}})}\;(\textsc{Task})
\qquad
\frac{\Gamma\vdash T:(A,s)\Rightarrow(B,s')\quad A\subseteq A'}
     {\Gamma\vdash T:(A',s)\Rightarrow(B,s')}\;(\textsc{Sub})
\\[2.2ex]
\frac{\Gamma\vdash T_1:(A,s)\Rightarrow(B,s')\quad\Gamma\vdash T_2:(B,s')\Rightarrow(C,s'')}
     {\Gamma\vdash T_1;T_2:(A,s)\Rightarrow(C,s'')}\;(\textsc{Seq})
\\[2.2ex]
\frac{\Gamma\vdash T_1:(A,s)\Rightarrow(B,s')\quad\Gamma\vdash T_2:(A,s)\Rightarrow(B,s')}
     {\Gamma\vdash \mathtt{Choice}(T_1,T_2):(A,s)\Rightarrow(B,s')}\;(\textsc{Choice})
\end{gather*}
The width-subsumption rule (\textsc{Sub}) widens the \emph{available} input fields---a task
needing $A$ still runs when more fields $A'\supseteq A$ are present---and leaves the business
state untouched. It is what lets the field component be matched by \emph{containment} while the
typestate component is matched by \emph{equality}, mirroring the two checks below; without it the
\textsc{Seq} rule would force the middle field-record to match exactly, which the tool does not
require.

The judgment is read over the workflow's \emph{task-adjacency} structure: \tool{} projects the
control-flow graph onto its \texttt{Task} states, following \emph{normal} edges through the
\emph{transparent} states (\texttt{Choice}, \texttt{Pass}, \texttt{Wait}) that carry no
contract, so $T_1;T_2$ denotes a task immediately followed---through any transparent
states---by another and $\mathtt{Choice}(T_1,T_2)$ a branch whose arms rejoin. \texttt{Catch}
(error) edges are \emph{not} projected, and \texttt{Map}/\texttt{Parallel} states are opaque
adjacency endpoints whose sub-machines are checked separately, so the discipline governs the
normal-flow task graph only. A declared \emph{protocol} $\delta\subseteq S_{\mathrm B}\times
S_{\mathrm B}$ lists the legal business-state transitions, and
$\mathrm{Valid}(W)\iff\forall T\in W.\,(s_{\mathrm{in}}^{T},s_{\mathrm{out}}^{T})\in\delta$.

The checks are edge-local. For each task-adjacency $T_1\!\to\!T_2$ whose endpoints both carry
the relevant declarations: the typed-contract check \code{SC1010} fires iff some required field
of $T_2$ is missing from $T_1$'s produced fields
($\mathit{in}(T_2)\not\subseteq\mathit{out}(T_1)$); the typestate check \code{SC2001} fires iff
$s_{\mathrm{out}}^{T_1}\neq s_{\mathrm{in}}^{T_2}$; and \code{SC2002} fires iff
$(s_{\mathrm{in}}^{T_2},s_{\mathrm{out}}^{T_2})\notin\delta$ when a non-empty protocol is
declared. A \texttt{Choice} imposes the predecessor's premise on each branch head. An edge
missing an endpoint declaration is passed over---the checks are vacuous there---so the guarantee
below is stated for \emph{fully declared} workflows.

\begin{theorem}[Declared-tier consistency]\label{thm:typecheck}
Let $W$ be \emph{fully declared}: every \texttt{Task} carries a contract and both typestate
labels, and a protocol $\delta$ is declared. Then \code{SC1010}, \code{SC2001}, and \code{SC2002}
emit no finding \emph{iff}
\begin{equation}
\begin{aligned}
  &\mathit{in}(T_2)\subseteq\mathit{out}(T_1)\ \text{and}\ s_{\mathrm{out}}^{T_1}=s_{\mathrm{in}}^{T_2}
  \ \text{for every adjacency } T_1\!\to\!T_2,\\
  &\text{and}\ (s_{\mathrm{in}}^{T},s_{\mathrm{out}}^{T})\in\delta \ \text{for every task } T.
\end{aligned}
\tag{$\ast$}\label{eq:edgecond}
\end{equation}
The three checks are therefore a \emph{sound and complete} decision procedure for
\eqref{eq:edgecond}: every $W$ satisfying \eqref{eq:edgecond} is accepted, and every violation is
reported by a finding naming the offending edge (\code{SC1010}/\code{SC2001}) or task
(\code{SC2002}). On the \emph{series-parallel fragment}---acyclic task graphs built from
sequencing and fully rejoining binary \texttt{Choice} whose arms share a common branch-point and
rejoin pair type---condition \eqref{eq:edgecond} is moreover equivalent to derivability
$\Gamma\vdash W:(A,s)\Rightarrow(B,s')$ together with $\mathrm{Valid}(W)$, so on that fragment the
checks decide well-typedness in the calculus above.
\end{theorem}
\begin{proof}
The equivalence ``no finding $\iff$ \eqref{eq:edgecond}'' is immediate: each check is exactly the
negation of one clause of \eqref{eq:edgecond}, evaluated independently per edge (per task for
\code{SC2002}), and full declaration leaves no edge vacuous. This holds for an arbitrary
normal-flow task graph, cyclic ones included---a back-edge is simply another adjacency subject to
\eqref{eq:edgecond}.

For the series-parallel fragment we show \eqref{eq:edgecond}$\iff\Gamma\vdash
W:(A,s)\Rightarrow(B,s')\wedge\mathrm{Valid}(W)$. ($\Leftarrow$) A derivation types each adjacency
by \textsc{Seq}, whose shared middle type $(B,s')$ forces $\mathit{in}(T_2)\subseteq\mathit{out}(T_1)$
on the field component (a preceding \textsc{Sub} step only enlarges the available fields) and
$s_{\mathrm{out}}^{T_1}=s_{\mathrm{in}}^{T_2}$ on the typestate component; \textsc{Choice} imposes
the same shared type on both arms, extending the equalities across a branch and its rejoin.
$\mathrm{Valid}(W)$ supplies the protocol clause, giving \eqref{eq:edgecond}. ($\Rightarrow$)
Assume \eqref{eq:edgecond}. Build the derivation bottom-up: \textsc{Task} gives each task
$T:(\mathit{in}(T),s_{\mathrm{in}}^T)\Rightarrow(\mathit{out}(T),s_{\mathrm{out}}^T)$. For a
sequential adjacency $T_1;T_2$ take the middle type
$(B,s')=(\mathit{out}(T_1),s_{\mathrm{out}}^{T_1})$: by \eqref{eq:edgecond},
$\mathit{in}(T_2)\subseteq\mathit{out}(T_1)=B$ and $s_{\mathrm{in}}^{T_2}=s_{\mathrm{out}}^{T_1}=s'$,
so \textsc{Sub} lifts the \textsc{Task} axiom of $T_2$ to $(B,s')\Rightarrow(\mathit{out}(T_2),
s_{\mathrm{out}}^{T_2})$ and \textsc{Seq} applies. A fragment \texttt{Choice} types both arms at
the common branch-point type $(A,s)$ and common rejoin type $(B,s')$ (each arm as above), so
\textsc{Choice} applies. The result derives $\Gamma\vdash W$, and the per-task protocol
membership gives $\mathrm{Valid}(W)$. Outside the fragment---non-rejoining or multiway branches,
type-misaligned rejoins, or cycles---no finite derivation in this three-constructor calculus need
exist even when \eqref{eq:edgecond} holds; there the edge-condition statement is the operative
one, and it is exactly what the tool decides.

Both readings are \emph{relative to the declared annotations}, and here that scoping is exact:
contracts and typestate labels are never inferred---only the idempotence/persistence/effect
annotations of Section~\ref{sec:tr:saga} are---so on this tier $\Gamma$ \emph{is} the declared
data and the three checks are unconditionally error-level. A workflow lacking a declaration is
outside the theorem's scope (the corresponding check is vacuous on that edge), not silently
downgraded to a warning.
\end{proof}

\section{Downstream-aware Saga compensation}
\label{sec:tr:saga}

\subsection{Model of failure and the post-commit region}

Let $W$ declare which tasks are \emph{persistent} and, for each persistent $P$, a declared
\emph{compensator} $C$ (the Saga pattern~\cite{garciaMolina1987}). Every \emph{fallible} state---a \texttt{Task}, \texttt{Map}, or
\texttt{Parallel}---may raise the generic error \texttt{States.TaskFailed}; an error caught by a
\emph{broad} \texttt{Catch} (one whose \texttt{ErrorEquals} contains \texttt{States.ALL} or
\texttt{States.TaskFailed}) is routed to the catch target, otherwise it propagates to workflow
failure. \texttt{Parallel} and \texttt{Map} are fallible black boxes; a nested persistent task is
analysed within its own sub-machine by the same procedure.

The completeness half of the guarantee below is stated relative to an explicit model of
execution, made precise here so the theorem is neither over- nor under-claimed.

\begin{assumption}[Failure model]\label{asm:saga-model}
\emph{(A1) Single post-commit fault}: on each execution we reason about the \emph{first} failure
after $P$ commits, and the compensation flow it triggers is itself fault-free---a standard
single-fault Saga assumption. \emph{(A2) Path feasibility}: \texttt{Choice} guards are abstracted,
so every normal edge is feasible (the analysis is path-insensitive). \emph{(A3) Unique generic
catcher}: no fallible state has two \texttt{Catch} clauses both matching \texttt{States.TaskFailed},
so the existential test over catchers coincides with \asl's first-match resolution for the modeled
error. \emph{(A4) Linear handlers}: from a broad \texttt{Catch} target, control follows a unique
normal successor until it reaches $C$ (as the typed DSL builds compensation chains), so
``reaches~$C$'' means ``executes~$C$''.
\end{assumption}

\begin{definition}[Post-commit region]
$R(P)$ is the set of states reachable from $P$'s normal successors by following \emph{normal}
(non-error) control-flow edges \emph{without passing through} $C$. These are exactly the states
that execute after $P$ has committed its durable effect and before any compensation.
\end{definition}

\begin{definition}[Saga-completeness]
A persistent task $P$ with compensator $C$ is \emph{Saga-complete} when every execution in which
$P$ commits its effect and a subsequent \emph{post-commit failure} occurs---a fallible state in
$R(P)$ raising the generic error, or a \texttt{Fail} terminal reached on a normal path in
$R(P)$---also executes $C$.
\end{definition}

\subsection{The \code{SC4011} algorithm}

\begin{lstlisting}[style=pseudo,caption={Downstream-aware Saga check (per persistent task).},label={lst:sc4011}]
SC4011(W, P, C):
  region = post_commit_region(W, P, C)         # the region R(P)
  foreach Q in region:
    if is_fallible(Q) and not failure_reaches(Q, C):
        fire SC4011 at Q                         # a post-commit failure at Q escapes C
    if is_fail_terminal(Q):
        fire SC4011 at Q                         # a normal path commits P then ends in failure

post_commit_region(W, P, C):                     # normal-flow closure of P's successors,
  return forward_closure(normal_successors(P),   #   not passing through C
                         edges = normal, without = {C})

failure_reaches(Q, C):
  # Q's own generic failure is caught broadly and the handler reaches C by normal flow
  return exists c in catches(Q):
           broad(c.error_equals) and reaches(c.target, C, edges = normal)

is_fallible(Q):      return Q.kind in {Task, Map, Parallel}
is_fail_terminal(Q): return Q.kind == Fail
broad(errs):         return "States.ALL" in errs or "States.TaskFailed" in errs
reaches(s, C, e):    return C in forward_closure({s}, edges = e)
\end{lstlisting}

\begin{theorem}[Compensation completeness under the declared model]\label{thm:saga}
Fix a workflow $W$ with \emph{declared} persistence and compensators, and adopt
Assumption~\ref{asm:saga-model}. For a persistent task $P$ with compensator $C$, \code{SC4011}
(Listing~\ref{lst:sc4011}) emits no finding \emph{iff} $P$ is Saga-complete---iff every fallible
$Q\in R(P)$ has a broad \texttt{Catch} whose target reaches $C$, and no normal path from $P$
reaches a \texttt{Fail} terminal in $R(P)$. Under the model, \code{SC4011} is a \emph{sound and
complete} decision procedure for Saga-completeness.
\end{theorem}
\begin{proof}
\emph{Soundness} (every finding is real, relative to A2). Suppose \code{SC4011} fires. In the
first case it fires at a fallible $Q\in R(P)$ with $\neg\,\mathit{failure\_reaches}(Q,C)$. By A2
the normal path $P\to\cdots\to Q$ is feasible, so some execution runs $P$ to commitment and
reaches $Q$; being fallible, $Q$ raises \texttt{States.TaskFailed} (A1: the single post-commit
fault). By A3 the catcher the test inspects is the one \asl{} would match, and it is absent,
non-broad, or routes to a handler not reaching $C$; by A4 ``reaches~$C$'' is exact, so $C$ does
not execute. This is a concrete post-commit failure that leaks $P$'s effect. In the second case
\code{SC4011} fires at a \texttt{Fail} terminal $Q\in R(P)$: the feasible normal path
$P\to\cdots\to Q$ commits $P$ and then fails with no compensation---a more direct witness.

\emph{Completeness} (every defect is found). Suppose $P$ is \emph{not} Saga-complete: some
execution commits $P$ and then a post-commit failure occurs without $C$. By A1 consider that
first failure. If it is a fallible state $Q$, then $Q\in R(P)$; by A3 its generic error resolves
to the single catcher the test inspects, and since $C$ did not run, that catcher does not route
to $C$ (A4 makes ``reaches'' exact), so $\neg\,\mathit{failure\_reaches}(Q,C)$ and \code{SC4011}
fires at $Q$. If instead the execution reaches a \texttt{Fail} terminal on a normal path without
$C$, that terminal lies in $R(P)$ and \code{SC4011} fires there. Conversely, if every fallible
$Q\in R(P)$ has its (unique, by A3) broad catcher reaching $C$ and no normal path bypasses $C$ to
a \texttt{Fail}, then by A4 every post-commit fault routes into a chain that runs $C$, so $P$ is
Saga-complete. Hence \code{SC4011} fires exactly on the non-Saga-complete workflows.

The guarantee is \emph{relative to the declared facts}: a compensator is never \emph{inferred},
so under pure inference the premise is absent and \code{SC4011} stays silent (in contrast to the
name-based, warning-level \code{SC4010}).
\end{proof}

\begin{remark}[Where the assumptions bite]\label{rmk:saga-limits}
Relaxing Assumption~\ref{asm:saga-model} costs completeness, not soundness. Without A1, a fault
\emph{inside} the compensation chain (e.g.\ a fallible compensator) can leak an earlier effect,
which \code{SC4011} does not model. Without A3, a fallible state whose first catcher matches
\texttt{States.TaskFailed} but routes away from $C$, shadowing a later broad catcher, is missed.
Without A4, a broad catcher whose handler \emph{may} reach $C$ on one branch and a \texttt{Fail}
on another passes the existential test yet can skip $C$; deciding this exactly needs $C$ to
\emph{postdominate} the handler, which the check approximates by reachability. Non-generic errors
(\texttt{States.Timeout}, service faults) caught by narrow catchers, and cross-branch aborts of a
\texttt{Parallel}/\texttt{Map} that strand a sibling's committed effect, are likewise outside the
model. Closing these is future work; the corpus in the main paper satisfies A1--A4.
\end{remark}

\subsection{Reverse-order chains}

Saga-completeness of a single $(P,C)$ pair is local, but a multi-step Saga must unwind \emph{all}
prior effects on a late failure. Take persistent tasks $P_1,\dots,P_n$ committed in order with
compensators $C_1,\dots,C_n$ wired as a \emph{reverse-order chain}: each $C_i$'s \emph{unique}
normal successor is $C_{i-1}$, and $C_1$'s is the failure terminal. Under this linear-chain
hypothesis---exactly A4 of Assumption~\ref{asm:saga-model} applied to the chain---running $C_i$
runs the whole prefix $C_i,C_{i-1},\dots,C_1$ in order, so applying the single-pair check of
Theorem~\ref{thm:saga} to each $(P_i,C_i)$ certifies that any post-commit failure unwinds the
full committed prefix. The hypothesis is essential: without the unique-successor wiring, per-pair
reachability certifies neither the \emph{order} nor the \emph{completeness} of the unwind---a
branch in the chain could skip some $C_j$ while still reaching a later $C_k$. This is how the
typed DSL builds a Saga by construction, which is why a correctly built DSL Saga raises no
\code{SC4011} and a one-edge break (redirecting one compensator to the failure terminal) makes
\code{SC4011} fire on precisely the upstream steps whose effects would then leak (the chain-break
experiment in the main paper's evaluation).

\section{Conservative analyses: concurrency and temporal (notes)}
\label{sec:tr:cons}

The concurrency and temporal analyses are native but deliberately \emph{conservative}: they
under-approximate to avoid false alarms and make no soundness claim. All but one of their
findings are warnings; the sole exception is the definite mis-configuration \code{SC6003}
(below), reported as an error.

\paragraph{Concurrency (\code{SC5001}/\code{SC5002}).}
\texttt{Parallel} branches and concurrent \texttt{Map} iterations each receive their own copy
of the document, so the document cannot race; only \emph{external} resources can. The check
mines each task's resource key---service plus the named target: \texttt{TableName},
\texttt{QueueUrl}, \texttt{TopicArn}, \texttt{Bucket}, \texttt{EventBusName}, ECS cluster/task
identifiers, Bedrock \texttt{ModelId}, and nested-workflow \texttt{StateMachineArn}/\texttt{Name}
pairs (statically named \texttt{states:startExecution} calls that provably share a child workflow
and execution name also collide)---and considers only \emph{overwriting}
writes---\texttt{PutItem}/\texttt{PutObject}, which are last-writer-wins. For child workflows, the
CloudFormation/SAM front end links sibling state machines through logical ids, \texttt{Ref}/\texttt{GetAtt},
\texttt{DefinitionSubstitutions}, and \texttt{StateMachineName}-derived ARNs, so the child's own write
set is recursively flattened when the definition is present in the artifact; externally deployed
children with no local definition remain a scope boundary. A task counts as a
write only when it is annotated non-idempotent (or persistent): an unannotated \texttt{PutItem}
is not flagged, a deliberate precision trade that keeps the analysis silent absent an
effect declaration and contributes to its zero-false-positive record. The check also
excludes \texttt{UpdateItem} (an atomic merge that commutes for disjoint attributes) and
fire-and-forget \texttt{Publish}/\texttt{SendMessage} (independent messages). \code{SC5001}
flags two \texttt{Parallel} branches that overwrite the same resource; \code{SC5002} a
concurrent \texttt{Map} (\texttt{MaxConcurrency}\,$\neq 1$) whose iterations overwrite a shared
resource while threading \emph{no} per-item data (no \texttt{.\$} reference in the write, so
every iteration clobbers the same item). Restricting to overwriting writes with a provably
shared key is what makes the analysis report a conflict only when two branches \emph{provably}
touch the same key---hence zero in-the-wild false positives, at the cost of not claiming
completeness.

\paragraph{Temporal (\code{SC6001}--\code{SC6003}).}
\code{SC6001} warns when a callback (\texttt{waitForTaskToken}) or Activity task has neither
\texttt{TimeoutSeconds} nor \texttt{HeartbeatSeconds}, so it can wait up to the platform's
one-year maximum. \code{SC6002} warns when a task's worst-case retry budget---the geometric
back-off sum $\sum_{k=0}^{a-1}\mathit{interval}\cdot\mathit{backoff}^{\,k}$ over $a$ attempts
(defaults $\mathit{interval}=1\mathrm{s}$, $\mathit{backoff}=2$, $a=3$)---exceeds a declared
machine \texttt{TimeoutSeconds}. \code{SC6003} is a definite mis-configuration:
\texttt{HeartbeatSeconds}\,$\geq$\,\texttt{TimeoutSeconds}, which the platform itself rejects
only at deploy time. \code{SC6001} is a warning because a missing timeout is occasionally
acceptable for a short-lived task; \code{SC6003} is the one error-level temporal finding, as
it is unambiguous.

\begin{remark}[Residual opacity beyond the modeled fragments]
\tool{} tracks workflow variables and a JSONata fragment---direct
\texttt{\$states}/\texttt{\$}\emph{name} references and object literals---as first-class provenance,
on equal footing with JSONPath: the document and variable shapes share one lattice and one fixpoint
(\S\ref{sec:tr:env}). What remains lifted to $\top$---soundly, at a cost in precision---is the
genuinely opaque residue: JSONata functions, filters, and arithmetic (\texttt{\$exists}, path
predicates, \texttt{\$merge}, string and math operators), the \texttt{\$states.context} object, and
the opaque results of external tasks. A richer JSONata value abstraction (for instance, tracking the
keys \texttt{\$merge} preserves) would recover some of this precision under the same lattice and
fixpoint; it is the natural next step, and the shape machinery already accommodates it.
\end{remark}

\section{Artifact and implementation notes}
\label{sec:tr:impl}

This section records the implementation and reproduction detail deferred from the main paper's
implementation and artifact sections. The artifact contains the verifier, the corpora used in the
paper, the mutation engine, the concrete execution oracle, the validator-panel scripts, the timing
harness, the real-bug replay harness, and the scripts that regenerate the main paper's figures and
tables. The AWS corpus has 193 Step Functions workflows; the cross-format corpus has 66 CNCF
Serverless Workflow examples.

Excluding tests, \tool{} is 4{,}552 lines of Rust, divided as follows: the workflow IR,
323 lines; the three frontends (\asl{} parse and emit, the typed DSL, and CNCF~\cite{cncf}), 784; annotation
resolution and inference, 275; the eight analysis passes and their manager, 1{,}730;
diagnostics, 115; the execution oracle, 269; the mutation engine, 316; and the command-line
interface together with the statistics, evaluation, and oracle harnesses, 740; a further 383
lines of in-tree tests are not counted above. The data-flow pass is by far the largest analysis
at 584 lines---the shape lattice of Section~\ref{sec:tr:df}, its transfer functions, and the
control-flow fixpoint---while the other seven passes range from roughly 85 to 240 lines each.
(These counts are for the artifact revision and drift with continued development.)

\paragraph{Packaging and distribution.}
The artifact ships a continuous-integration pipeline that builds and tests the binary on Linux,
macOS, and Windows on every change, and a release pipeline that packages it for language package
managers---a \texttt{crates.io} crate (\texttt{cargo install stepcheck}) and an \texttt{npm}
wrapper package that fetches the platform binary---so a Rust- or Node-centric CI adopts it with
one command; a copy-pasteable GitHub Actions gate is included. With \texttt{--json}, findings
emit one stable object per finding carrying a code, severity, offending state, message, and
remediation note (e.g.\ \code{SC1101} at \code{ChargeCard}: ``reads \texttt{\$.total}, a field no
execution reaching it can have produced,'' noting that no state on any path binds it).

The command-line tool exposes \texttt{check}, \texttt{infer}, \texttt{scan}, \texttt{stats},
\texttt{eval}, \texttt{oracle}, \texttt{eval-pairs}, \texttt{mutate}, \texttt{emit}, and
\texttt{demo}. \texttt{check} runs the eight passes; \texttt{emit} round-trips a workflow
through the \asl{} emitter unchanged, the control form against which the multi-validator
baseline of the main paper scores each injected defect; \texttt{oracle} runs the data-flow
soundness cross-check; \texttt{eval-pairs} drives the real-bug benchmark by replaying each
pre-fix and post-fix workflow; \texttt{mutate} injects one defect of a chosen class; and
\texttt{demo} emits a built-in typed-DSL example to \asl.

The data-flow commands also expose the \texttt{--result-shapes} ablation used in the main paper's
evaluation. This flag leaves the default opaque-task transfer unchanged but, for the ablation run,
treats declared task-output schemas and a small table of AWS service-result envelopes as closed
records. On the strict-input missing-field mutation study, this improves \code{SC1101} recall from
84/126 to 88/126; the paired bounded-loop oracle run confirms all 92 \code{SC1101} reports emitted
under the ablation with zero counterexamples.

The evaluation is reproducible offline on a commodity laptop. A Rust~1.96 toolchain builds the
verifier and runs the mutation study, data-flow oracle, timing experiments, and real-bug replay.
Node.js and Ruby drive the two local third-party validators used in the baseline
(\texttt{asl-validator} and \texttt{statelint}), and Python renders the figures. No AWS account or
credentials are required to reproduce the paper's tables, figures, mutation results, data-flow
oracle, timing results, or local validator baselines. AWS's server-side
\texttt{ValidateStateMachineDefinition} check is opt-in (\texttt{STEPCHECK\_AWS=1}), creates no
resources, and has recorded outputs committed with the artifact so its numbers reproduce offline.
The only steps that touch an AWS account are the optional deployment checks reported in
the main paper (Express, Standard, and live-callback runs; scripts
\texttt{infra/deploy\_run\{,\_standard,\_callback\}.sh}); no other result depends on them, and the
captured execution evidence is committed with the artifact.

The \asl{} emitter is currently lossy for a workflow's pure I/O-processing fields: among others
it does not re-emit \texttt{InputPath}, \texttt{OutputPath}, or \texttt{ResultSelector}, nor the
machine-level \texttt{TimeoutSeconds}; it also drops the \texttt{QueryLanguage} directive (a
JSONata machine re-emits without it) and canonicalises a \texttt{Map}'s \texttt{ItemProcessor} to
the older \texttt{Iterator} form. The AWS round-trip is therefore exercised on DSL-authored
workflows that use none of these fields, where the deployed state machine is behaviourally
equivalent to the verified IR.

\section{Extended limitations}
\label{sec:tr:limits}

This section expands the threats-to-validity discussion of the main paper's evaluation.

\paragraph{Corpus composition.}
The corpus comes from public AWS and CNCF sample repositories and skews small (median 6 states,
maximum 33). We do not include proprietary industrial definitions because workflow files are
sensitive artifacts: they expose resource names, account topology, and business-process
structure, and no anonymization we are aware of preserves the semantic features the analyses
consume (JSONPath references, resource keys, task names for inference). The generated chain,
fan-out, and nested-\texttt{Map} sweeps exercise scale and topology, but they cannot stand in
for industrial semantic diversity---richer contracts, deeper Saga chains, and organization-specific
naming conventions that inference has not been tuned on.

\paragraph{Deployment check.}
The AWS deployment check exercises three execution models on StepCheck-verified definitions, each
deployed unmodified (only \texttt{FunctionName} bound to a stand-in echo Lambda) and each reaching
\texttt{SUCCEEDED} before teardown: (i)~the DSL-authored order workflow as a synchronous
\emph{Express} machine; (ii)~the same definition as a \emph{Standard} machine executed
asynchronously, whose durable history retains all 24 events (Express retains none); and (iii)~a
DSL-verified order variant with a bounded \texttt{.waitForTaskToken} approval gate deployed as
Standard (Express cannot run this integration), where the execution pauses at the callback, an
external \texttt{SendTaskSuccess} resumes it, and it completes in 30 events. The callback carries a
\texttt{TimeoutSeconds} bound, so the temporal pass raises no \code{SC6001}; removing the bound makes
StepCheck flag it before deployment, tying the check to the deployed behaviour. To confirm the
verified artifact runs correctly beyond a single execution, we additionally run the order workflow
\textbf{100} times in each non-callback mode: all complete successfully (Express 100/100, Standard
100/100), with per-execution wall time of p50/p95/p99 $\approx$72/130/165\,ms (Express, mean
$\approx$84\,ms) and $\approx$345/456/484\,ms (Standard, mean $\approx$361\,ms, five state
transitions per run). We report these distributions as evidence of stable completion across
execution models, not as a cross-mode AWS performance benchmark. StepCheck itself is also wired as a
live gate: \texttt{.github/workflows/gate.yml} runs the blocking check (\texttt{-{}-deny-warnings})
on every push/PR over the deployment artifacts and includes a negative test that injects a defect and
fails the job unless the gate blocks it. What the deployment check still does \emph{not} exercise is
Activity workers, service-integration failures, or a broad multi-workflow performance benchmark;
those behaviours remain covered by the static analyses and the offline oracle rather than by live
execution.

\paragraph{Semantic drift.}
Declared premises are checked for syntactic rot (\code{SC0011}/\code{SC0012}) but not for truth:
a sidecar that says a Lambda is idempotent when its implementation is not will pass the
declared-tier checks. Generating the sidecar from code or infrastructure-as-code---the
generated-metadata premise source---is the deployment pattern that keeps declared facts in sync;
inference is the fallback and always stays at warning level.

\begin{thebibliography}{9}
\bibitem{stepcheck} Anonymous Author(s). Static Verification of AWS Step Functions with Sound
Data-Flow Analysis and Workflow Semantics. Under submission to ICSOC, 2026. (The main paper this
report accompanies.)
\bibitem{cousot1977} P.~Cousot and R.~Cousot. Abstract interpretation: a unified lattice
model for static analysis of programs by construction or approximation of fixpoints.
In \emph{POPL}, 1977.
\bibitem{garciaMolina1987} H.~Garcia-Molina and K.~Salem. Sagas.
In \emph{ACM SIGMOD}, 1987.
\bibitem{stromYemini} R.~E.~Strom and S.~Yemini. Typestate: a programming language concept for
enhancing software reliability. \emph{IEEE Transactions on Software Engineering}, 12(1), 1986.
\bibitem{delineFahndrich} R.~DeLine and M.~F\"ahndrich. Typestates for objects.
In \emph{ECOOP}, 2004.
\bibitem{jsonpath} S.~G\"ossner, G.~Normington, and C.~Bormann (eds.).
\emph{RFC 9535: JSONPath --- Query Expressions for JSON}. IETF, 2024.
\bibitem{cncf} Cloud Native Computing Foundation. \emph{Serverless Workflow Specification}.
\url{https://serverlessworkflow.io/}.
\bibitem{awsasl} Amazon Web Services. \emph{AWS Step Functions Developer Guide / Amazon
States Language specification}. \url{https://states-language.net/spec.html}.
\end{thebibliography}

\end{document}
